\section{第~\ref{sec:serocomputer}~章：\nt{SERO}神经元}

\nt{SERO}神经元本身的性质（节律、自抑制、由受体决定的符号、子系统）是直接测量的；本章的计算——调节器、滤波器、逻辑——由本章的方程推出；$\tau_c$、血清素的扩散系数和释放位点数是估计；回路在脑中充当电平调节器是假说（在中缝核内自抑制是点对点的，只造成短暂停顿），而\nt{SERO}充当视野则是有行为实验支持的假说。

{\footnotesize
\setlength{\tabcolsep}{3pt}\renewcommand{\arraystretch}{1.15}
\begin{longtable}{>{\raggedright}p{0.5cm}>{\raggedright}p{4.0cm}>{\raggedright}p{5.6cm}>{\raggedright}p{2.3cm}>{\raggedright\arraybackslash}p{1.7cm}}
\toprule
序 & 命题 & 实验与测量内容 & 来源 & 状态 \\
\midrule\endfirsthead
\toprule
序 & 命题 & 实验与测量内容 & 来源 & 状态 \\
\midrule\endhead
1 & 血清素作用的符号由接收者的受体选择（命题~\ref{prop:receptor}） & 血清素受体按药理和机制分为若干家族：5-HT$_{1A}$经 G 蛋白打开钾通道并抑制细胞，5-HT$_{2A}$和离子型的 5-HT$_3$则起兴奋作用。同一种递质在不同细胞中产生相反的响应。 & \cite{BarnesSharp1999} & 已测量 \\
2 & 清醒时\nt{SERO}神经元平稳地每秒发放几个脉冲 & 记录自由活动猫的中缝背核神经元：安静清醒时为缓慢而有节律的放电，$2.8\pm0.2$次/秒；慢波睡眠中频率下降$34$--$68\,\%$，快速眼动睡眠中下降$98\,\%$。麻醉大鼠中为$1.7\pm0.5$次/秒，非常规则。 & \cite{TrulsonJacobs1979, GagnonParent2014, JacobsAzmitia1992} & 已测量 \\
3 & \nt{SERO}神经元是起搏器：不需要每个脉冲都有一次推动，但需要平稳的背景输入（脑片中是经$\alpha_1$受体的去甲肾上腺素，体内来自\nt{NORA}） & 在脑片中细胞缓慢而平稳地放电，每个脉冲之后有持续$200$--$800\ms$的后超极化。脑片中自发工作的细胞少于体内：平稳节律需要经$\alpha_1$受体的去甲肾上腺素紧张性输入，阻断这些受体则节律熄灭。 & \cite{VandermaelenAghajanian1983} & 已测量 \\
4 & 受体几乎都是代谢型的；响应缓慢：在大约一秒内升起又衰减 & 除 5-HT$_3$外，所有受体家族都与 G 蛋白偶联。在脑片中，诱发的血清素释放经 5-HT$_{1A}$产生抑制性电流，在一秒内升起又衰减。 & \cite{BarnesSharp1999, CourtneyFord2016} & 已测量 \\
5 & 在目标区域血清素释放到体积中，而不只是释放到间隙中；在中缝核内部，对邻居的抑制是点对点的 & 脑片中的快速循环伏安法：每个脉冲的释放量在有突触的区域和几乎没有突触的中缝核中相同；不受摄取阻断和受体阻断的影响；每个末梢的分子数远少于转运体和受体的数目，而细胞外浓度与主要受体的亲和力相符。在中缝核内，经 5-HT$_{1A}$的抑制是点对点的：转运体不让血清素在间隙外积累。 & \cite{Bunin1998, CourtneyFord2016} & 已测量 \\
6 & 按~\eqref{eq:seroneuron}，浓度因泵回而下降；$\tau_c$量级为一秒是估计 & 活体脑中的伏安法：细胞外血清素主要受再摄取和降解限制，而不是受合成限制。所引文献中没有$\tau_c$的直接测量；其量级是本章的估计。 & \cite{Hashemi2012serotonin} & 已测量；$\tau_c$为估计 \\
7 & 一个起搏器实现“非”和“或非”；\nt{SERO}逻辑的完备性（命题~\ref{prop:serocomplete}，图~\ref{fig:serogates}） & 由~\eqref{eq:seroneuron}推出：受抑制的起搏器是“或非”，而“或非”在功能上是完备的。没有用\nt{SERO}神经元搭建逻辑的实验；在脑中，体积彼此分离的条件并不成立。 & 本章推导 & 理论 \\
8 & 血清素经\nt{SERO}神经元自身的受体抑制它 & 在麻醉大鼠中，静脉注射和离子电泳给予的 5-HT$_{1A}$激动剂抑制中缝背核神经元放电，其剂量—反应关系与血清素本身相同；5-HT$_{1B}$激动剂几乎无作用。在脑片中，邻近细胞的内源性血清素经 5-HT$_{1A}$产生电流和放电中的短暂停顿，不改变平均频率。 & \cite{Sprouse1987, CourtneyFord2016} & 已测量 \\
9 & 在模型中，经公共体积的回路是电平调节器：离散和时间常数缩小为$1/(1+G)$~\eqref{eq:feedback}；回路在脑中这样工作是假说 & 负反馈放大器的经典公式；本章重新推导。\nt{SERO}系统的环路增益$G$没有测量过。已测量的是：在脑片中自抑制是点对点的，造成短暂停顿而不改变平均频率。 & \cite{Black1934feedback, CourtneyFord2016} & 理论；在脑中为假说 \\
10 & 浓度是频率的滑动平均，即低通滤波器~\eqref{eq:lowpass}，图~\ref{fig:lowpass} & 线性方程~\eqref{eq:seroneuron}的解；图是按该公式在$\tau_c=1\,\text{s}$时的计算。\texttt{models/}中没有单独的模型。 & 本章推导 & 理论 \\
11 & 间隙的曲折使小分子的扩散减慢约$2.6$倍 & 点源法：离子电泳注入四甲基铵，在脑片和活体脑中用离子选择电极记录其浓度。曲折度$\lambda\approx1.6$，即有效$D$比自由扩散小$\lambda^2\approx2.6$倍；细胞间隙体积分数约为$0.2$。这些工作没有测量血清素本身在组织中的扩散系数；本章中它是估计。 & \cite{Sykova2008} & 已测量 \\
12 & 独立“导线”的长度$\ell\sim\sqrt{D\tau}\approx20\um$~\eqref{eq:diffusionlength}；导线数受这类区域的数目限制（命题~\ref{prop:volumewires}） & 扩散定律并代入$D$和$\tau$的估计（第~6、11~行）；命题~\ref{prop:volumewires}是推论。没有直接计数容积传递独立通道的实验。 & 本章推导 & 理论 \\
13 & 一个\nt{SERO}神经元的输出散布在脑的巨大区域上；释放位点估计约$\sim10^5$个 & 大鼠中缝背核单个神经元的完整重建：所有上行轴突都高度分支，一个细胞的轴突总长可达$18.7$~cm，一个细胞的分支可到达纹状体、前额叶皮层和杏仁核。每个细胞的释放位点数没有直接计数。 & \cite{GagnonParent2014} & 已测量；$10^5$为估计 \\
14 & \nt{SERO}神经元分为几个子系统，输出和响应各不相同 & 小鼠的病毒—遗传标记：投射到杏仁核和投射到额叶皮层的神经元在分支上几乎不重叠，接收不同的输入，对厌恶刺激的响应相反；激活或抑制每一组对行为的改变各不相同。 & \cite{Ren2018} & 已测量 \\
15 & 指数折扣下的耐心条件$D<\tau_\gamma\ln(R/r_0)$~\eqref{eq:patience}；在实测的双曲折扣下为$D<\tau_\gamma(R/r_0-1)$ & 由折扣$e^{-D/\tau_\gamma}$或$1/(1+D/\tau_\gamma)$推出；本章推导。猴子在秒级延迟上的选择：折扣更接近双曲线而非指数；\nt{DOPA}神经元对延迟奖励预示信号的响应随延迟下降，同样符合双曲线。 & 本章推导；\cite{KobayashiSchultz2008} & 理论 \\
16 & \nt{SERO}设定视野$\tau_\gamma$（第~\ref{sec:horizon}~节） & 在小鼠中用光遗传学激活中缝背核\nt{SERO}神经元，会延长对延迟奖励的等待，并提高在奖励变稀少时寻找奖励的坚持程度。在人类中，降低血清素（无色氨酸饮食）使延迟奖励的折扣更陡。浓度如何变成$\tau_\gamma$尚不清楚。 & \cite{Doya2002, Miyazaki2014, Lottem2018, Schweighofer2008} & 假说 \\
\bottomrule
\end{longtable}}
